\documentclass[10pt,a4paper]{article}
\usepackage[utf8]{inputenc}
\usepackage[T1]{fontenc}
\usepackage[margin=2.2cm]{geometry}
\usepackage{booktabs}
\usepackage[expansion=false]{microtype}
\usepackage[numbers,sort&compress]{natbib}
\usepackage[colorlinks=true,linkcolor=black,citecolor=black,urlcolor=black]{hyperref}
\usepackage{enumitem}
\setlist{nosep}
\title{Ikhtibar: a comparative evidence base on pre deployment testing of AI systems, with baselines, a verification audit and a checklist tool}
% Affiliation, acknowledgements and funding to be completed by the author before submission.
\author{Huda Mueen\\United Arab Emirates}
\date{Draft of 28 October 2026}
\begin{document}
\maketitle

\begin{abstract}
Governments, standards bodies and institutes have published many documents that mention testing AI systems before they are deployed, but the documents differ in force: some require testing, some recommend it and some only describe it. We assemble an evidence table of {{n_rows}} provisions from public primary documents, of which {{n_test_head}} are test provisions that pass a verification screen, drawn from {{bl_frameworks}} documents in {{bl_families}} families. Of these, {{head_mc}} ({{head_mc_pct}}) are mandatory or conditional, {{head_vg}} ({{head_vg_pct}}) are voluntary or guidance, and {{head_am}} ({{head_am_pct}}) are ambiguous. We test this finding and the tooling around it instead of only asserting it. Simple baselines predict the binding status of a provision well only when other provisions of the same document were seen in training; on held out families a wording rule does better than a learned text model and the document prior falls to the majority rate. Two blind second codings of {{ag_n}} rows agree closely with each other (Cohen's kappa {{ag_k_blind}} on five classes) and less with the original coding (kappa {{ag_k_orig_range}}). The finding that binding provisions are a minority {{rb_survives}} we tested, with a binding share of at most {{rb_max}}\% under the most favourable assumptions. A fresh independent check of a random sample found a strict error rate of {{vf_err_s_pct}}\% ({{vf_err_s}} of {{vf_read}} readable rows, 95\% interval {{vf_err_s_ci}}), but {{vf_unread_pct}}\% of the sample could not be read, including {{vf_eu_sampled}} sampled European Union rows of which only {{vf_eu_read}} could be read, so the rate is not representative of the rows behind the headline. A three valued checklist engine agreed with {{ev_acc}}\% of {{ev_dec}} preregistered item decisions (after two defects found by the scenarios were fixed) and never wrongly excluded an item. We state what the study does not show, list the open items, and release the data, code and tests.
\end{abstract}

\section{Introduction}
Organisations that build or buy AI systems are told, by many sources, to test them before use. A procurement officer reading these sources meets three different things under one word. A statute may oblige a named actor to test. A national institute may recommend a testing practice. A standard may describe what testing typically involves. Treating all three as requirements overstates the obligations on an organisation, and treating all three as advice understates them. Which of the three applies depends on the document and on who the organisation is.

This paper presents Ikhtibar (Arabic for test), a comparative evidence base on what published documents say about pre deployment testing of AI systems, coded by binding force and by the type of test. The work began as a gap analysis for public sector deployment in the United Arab Emirates, and it carries a section on what the retrieved UAE sources do and do not contain. The main contribution is not the table itself but the way we test it. A table of extracted provisions can look rigorous while resting on one analyst's coding and on rows that were never checked against the source. We therefore ask five questions of our own work:
\begin{enumerate}
\item Are the codes predictable from simpler signals, and which signal drives them? (baselines)
\item Do independent coders reproduce the codes? (coding reliability)
\item Does the headline finding survive reasonable changes to the coding choices? (robustness)
\item How often does a row misstate its source, estimated on a fresh random sample? (verification)
\item How much of the relevant literature does the table miss? (a completeness pilot)
\end{enumerate}
We also evaluate the checklist tool that turns the table into questions for an organisation, against scenarios written before the tool was run and against naive designs. Section~\ref{sec:limits} states what remains open. This is research analysis and not legal advice.

\section{Related work}
Comparative studies of AI regulation across jurisdictions exist, for example comparisons of the European Union, the United States and China \citep{chun2024comparative,du2025toward}, a seven category framework applied to five jurisdictions \citep{roberts2023comparative}, and a comparison of how risk is classified and enforced \citep{almaamari2025between}. These works compare regulatory approaches. We narrow the question to one activity, testing before deployment, and we record the force of each provision at the level of the clause, with a verification status per row.

Work on audits and evaluations explains what testing can and cannot establish. Internal auditing across the development lifecycle was proposed as a complement to external audits that happen after deployment \citep{raji2020closing}. Auditing at the level of providers, pre release models and applications was proposed for language models \citep{mokander2023auditing}, black box access was argued to be insufficient for rigorous audits \citep{casper2024blackbox}, and a survey of 390 audit tools and 35 practitioner interviews found that most tools support evaluation and standards management while support for harm discovery and audit communication is thin \citep{ojewale2024towards}. Evaluation methods in the research literature include holistic benchmarks \citep{liang2023holistic}, evaluations for dangerous capabilities \citep{shevlane2023model,phuong2024evaluating}, automated red teaming \citep{perez2022red}, and a sociotechnical view of safety evaluation \citep{weidinger2023sociotechnical}. A review of industry red teaming found that practices vary widely in goals and decisions informed \citep{feffer2024red}. These works describe how to test. We ask where documents say that testing must, should or does occur.

On the European Union Act, the force and mechanism of conformity assessment and post market monitoring have been analysed \citep{mokander2021conformity}, the use of harmonised standards has been examined \citep{gornet2024european}, and the reliance of the Act on self certification has been criticised \citep{wachter2024limitations}. Two of these analyses concern the 2021 proposal \citep{mokander2021conformity,veale2021demystifying}, while our table uses the amended text of 2026. Our table is consistent with the observation that the Act leaves much to the provider, but it covers the Act as one family among sixteen. For frameworks in other places, the voluntary testing framework of Singapore \citep{lee2023aiverify,harjani2025singapore} and the NIST risk management framework \citep{tabassi2023airmf} are examples of guidance that organisations may adopt. For the Gulf, a document analysis of the six national AI strategies of the GCC finds reliance on soft regulation through strategies and ethical principles, with no AI specific law in any member state \citep{albous2025gcc}, and a legal research overview of the UAE, Saudi Arabia and Bahrain describes non binding ethics guidelines and a non binding charter for the UAE \citep{hofverberg2024falqs}. A comparison of public sector AI adoption in the UAE and Kuwait studies institutional mechanisms rather than testing \citep{albous2025kuwaituae}. We found no study of testing requirements for AI in the UAE or the Gulf; this is a limit of our literature search, which ran under a restricted search budget, not a claim that none exists. For the public sector, impact assessment frameworks for agencies were proposed early \citep{reisman2018aia} and cautioned against organisation convenient metrics \citep{metcalf2021aias}; a framework for supreme audit institutions covers public sector AI across its lifecycle \citep{brand2025sai}, a national audit office found that only three of nine government algorithms met all of its requirements \citep{rekenkamer2022audit9}, a study of the transparency recording standard in policing reports fears of bureaucratic burden \citep{zilka2023police}, and a government roadmap sets out plans for an independent assurance market \citep{dsit2025assurance}. These works concern assessment and oversight of public sector systems; none compares pre deployment testing requirements across regulators, which is the gap this table addresses.

Our methods draw on three bodies of work. Language models are increasingly used for annotation, with results that range from agreement above crowd workers \citep{gilardi2023chatgpt,tornberg2023chatgpt} to reliability that varies with prompt wording and temperature \citep{reiss2023testing}, and with comparable performance to trained annotators on legal concepts when guidelines are supplied \citep{savelka2023can}; a survey covers the practice \citep{tan2024large} and a benchmark of thirteen models found fair agreement with human labellers \citep{ziems2024can}. Because our coding was assisted by a language model, these results motivate the blind second codings reported here. Agreement is reported with Cohen's kappa \citep{mchugh2012interrater} and Krippendorff's alpha \citep{krippendorff2011computing}. Kappa depends on the marginal distribution of labels, and sample size requirements for kappa tests depend on it \citep{bujang2017guidelines,wongpakaran2013comparison}, so we report percentage agreement beside kappa. Finally, language models are known to fabricate references \citep{fabrication2023chatgpt,chelli2024hallucination,athaluri2023exploring,naser2026cite}, which is a risk for any evidence table built with model assistance and a reason that rows here were checked against the source where the source could be read. Capture recapture methods have been used to estimate what a literature search misses \citep{spoor1996use,kastner2009capture,stelfox2013capture}, although their usefulness has been questioned for searches with duplicate identification problems \citep{motschall2009number}. Our completeness check is a simpler pilot against an independently built reference list.

\section{Evidence base and method}
\subsection{Frame, rows and coding}
We included primary documents from governments, national institutes and standards bodies that were retrieved and read, and that contain at least one provision on testing, evaluation, assessment, validation, conformity, red teaming, audit or verification of an AI system before or at deployment. Secondary sources were used only to locate primary ones. Each provision is a row with the clause or page, a paraphrase, the source address, the retrieval date, a binding status and a verification status. Of {{n_rows}} rows, {{n_test_all}} are test provisions (a provision that calls for, or describes as practice, an evaluation that produces evidence), {{n_enabler}} are enablers such as documentation or logging, and {{n_context}} are context such as dates and penalties.

Binding status has five values. \emph{Mandatory} means binding on the stated actor with the duty in operative text. \emph{Conditional} means binding only if a trigger applies, such as a risk class. \emph{Voluntary} is a formal scheme or standard an organisation may adopt. \emph{Guidance} is a recommendation or a description of practice. \emph{Ambiguous} is used where the force of the source cannot be settled from the text. Each test row also has a primary test type from a taxonomy of fifteen types (for example performance, security, red teaming, fairness, conformity assessment and post deployment monitoring). We report a headline over the {{n_test_head}} test rows that are not unverified, unverifiable or dependent on a portal summary; {{n_unreliable_rows}} rows are excluded for these reasons and their effect is bounded in Section~\ref{sec:robust}. Coding is judgemental and was assisted by language model agents, which we disclose in Section~\ref{sec:ethics}.

\subsection{Access rule}
Only standard fetch and search tools were used. If a site rejected a fetch, the attempt was recorded and no workaround was used (no changed user agent, cache, mirror or third party copy). The row is then marked as unverified because the source was blocked, it is excluded from the headline, and the document is listed for a person to download. {{n_blocked}} rows were blocked in this way at the end of the third phase. Two further kinds of row rest on weaker access: {{n_translated}} rows were read from an English translation or a machine rendering of the original (Chinese, Korean and Japanese sources, and some UAE documents), and {{n_iso}} rows ({{n_iso_head}} of them headline rows) come from ISO catalogue pages, where the paywalled standard text was not read. Each row records which applies.

\section{What the evidence says}
\textbf{Binding provisions are a minority.} Of {{n_test_head}} headline test provisions, {{head_mc}} ({{head_mc_pct}}) are mandatory or conditional, {{head_vg}} ({{head_vg_pct}}) are voluntary or guidance and {{head_am}} ({{head_am_pct}}) are ambiguous. Binding text is concentrated in a few families, led by the European Union Act, with others in health regulation in Abu Dhabi, China, the United States federal government (OMB), and the UAE and Korea.

\textbf{Interpretation: binding text seldom says how.} In the binding rows we read most closely, the text requires that testing happens and is documented, with the metrics and thresholds left to the provider. This is the author's reading of a few cited rows and not a counted result. The EU Act, for example, leaves the choice of metrics to the provider, and OMB requires pre deployment testing plans for high impact AI.

\textbf{Some test types are rarely binding.} Of {{n_RED}} headline rows coded red teaming or dangerous capability evaluation, only {{n_RED_binding}} are binding. Human oversight validation has {{n_OVER}} headline row and content safety evaluation has {{n_CSAFE}}, none binding, but these counts are too small to support more than a note. Shares differ by type (the full table is in the report), so counts of binding rows should not be ranked without the number of rows behind each.

\textbf{Required, recommended and described.} Among {{nf_n}} headline test rows that are not binding, two coders separated {{nf_obl}} that use obligation wording inside a non binding document, {{nf_rec}} that recommend a test and {{nf_des}} that only describe practice; {{nf_dis}} were disputed. The distinction matters because a reader who sees ``shall'' in a guidance document may take it for a requirement.

\textbf{UAE public sector.} We report what the retrieved sources contain and make no claim about what the UAE requires beyond them. One cross sector instrument that uses obligation wording for government entities, the national cyber security policy for AI, is about AI security. In the texts retrieved we found no binding cross sector requirement to test performance or fairness before deployment for non health uses, benchmarks and thresholds are not defined, independent assurance is contingent, and no row requires language specific performance testing. Many instruments could not be read or were not public. Absence in this table is not evidence of absence in law.

\section{Evaluating the evidence base}
\subsection{Baselines}\label{sec:baselines}
We ask how well the two main codings can be predicted by methods much simpler than reading the document, using the {{bl_n}} headline test rows. Methods were chosen by the analyst, who had seen the table, and committed with the code: a majority class; a document prior (the most common label among training rows of the same document); a wording rule for binding status (obligation words against recommendation words) and a keyword rule for test type built from the taxonomy definitions; and TF-IDF with logistic regression, with and without a document indicator. Protocols are {{bl_repeats}} times repeated stratified five fold cross validation (CV) and leave one family out (LOFO), which removes every document of a family from training. Table~\ref{tab:bl} gives accuracy; intervals and macro F1 are in the repository.

\begin{table}[t]\centering\small
{{tex_baselines}}
\caption{Accuracy (\%) of baselines for binding status in five classes, binding status in three classes and primary test type. n/a means the method does not apply to that task.}\label{tab:bl}
\end{table}

Binding status is largely a property of the document, in the sense that a document prior reaches {{bl_fw5}}\% when other provisions of the same document were seen in training, against {{bl_maj5}}\% for the majority class. This is close to tautological, since provisions of one document tend to share its force, and the prior falls back to the majority rate under LOFO ({{bl_fw5_lofo}}\%) because it has nothing to look up. On held out families the wording rule is the better method. A bag of words model reaches only {{bl_tfidf5}}\%. For test type, the keyword rule ({{bl_kw}}\%) beats the learned model ({{bl_tfidf_type}}\%) and the majority class ({{bl_maj_type}}\%); with a few hundred rows and fifteen classes a learned model has too little data. These scores show that the coding is regular, not that it is correct, since the labels are the analyst's own, and the paraphrase being classified was written by the same analyst. Intervals are bootstrap intervals over rows and ignore clustering by document, so they understate uncertainty; the document level bootstrap in Section~\ref{sec:robust} is wider. The keyword lexicon was written by an analyst who had seen the table. They also set the bar for any later classifier: it must beat the document prior and the keyword rule, not only the majority class.

\subsection{Coding reliability}
Two coders who saw the row paraphrase, the document name and the locator, and a written codebook (not the original codes), recoded {{ag_n}} randomly chosen headline test rows. Table~\ref{tab:ag} gives percentage agreement and Cohen's kappa for each pair, and Krippendorff's alpha over three coders.

\begin{table}[t]\centering\small
{{tex_agreement}}
\caption{Agreement (percent / kappa) between the original coding and two blind coders (B2, B3), and alpha over the three.}\label{tab:ag}
\end{table}

The blind coders agree closely with each other (kappa {{ag_k_blind}} on five classes). Each agrees less with the original coding (kappa {{ag_k_orig_range}}). In {{ag_odd_original}} of the {{ag_odd_total}} five class disagreements where exactly one coder differs (one more row has three different codes), the original is the odd one out. The original coder used ambiguous on {{ag_amb_orig}} of {{ag_n}} rows, against {{ag_amb_b2}} and {{ag_amb_b3}} for the blind coders, and the commonest pattern is ambiguous in the original against guidance in both blind codings. Two things could produce this and we cannot separate them: the original coder read the whole document and the blind coders read one provision, and the written instructions differed (the blind codebook asks for ambiguous to be used sparingly). Because the coders saw the analyst's paraphrase and the document name, agreement is probably higher than it would be on raw text. Agreement among coders is not accuracy, and all three coders are language models, so errors may be correlated. A human second coding of 100 rows is prepared and is the planned remedy.

\subsection{Robustness of the headline}\label{sec:robust}
Table~\ref{tab:rb} recomputes the binding share of test provisions under changes that a sceptical reader might propose: dropping each family in turn, adding back the {{n_test_excl}} excluded test rows as all binding, counting every ambiguous row as binding, doing both, weighting each document equally, and resampling documents.

\begin{table}[t]\centering\small
{{tex_robust}}
\caption{Binding (mandatory or conditional) share of test provisions under alternative assumptions.}\label{tab:rb}
\end{table}

The statement that binding provisions are a minority {{rb_survives}} we tried. The most favourable combination gives {{rb_max}}\%. The result is not driven by a few long documents: of {{rb_docs}} documents with headline test rows, {{rb_docs_any}} ({{rb_docs_any_pct}}\%) contain any binding test provision. The size of the minority does depend on coding choices, so it should be quoted with its range.

\subsection{Fresh independent verification}\label{sec:verify}
We drew {{vf_sampled}} rows at random with a fixed seed from the {{vf_pop}} rows that earlier checks had marked verified, and independent verifiers saw only the row and the rules and reread the primary source. Of these, {{vf_read}} could be read in the working session ({{vf_unread}} could not, {{vf_unread_pct}}\%, mostly because of a session level domain restriction and not because sites refused). Table~\ref{tab:vf} gives the error rates on readable rows. After strict grading {{vf_err_s}} rows had an error (they are listed in the repository). Accepted corrections were applied to the table. If every unreadable row were an error, the rate for the whole sample would be as high as {{vf_bound_hi}}\%, so we report the rate for readable rows only. This rate is for rows that earlier checks had already passed and for the readable part of the sample: only {{vf_eu_read}} of {{vf_eu_sampled}} sampled European Union rows could be read, and that family supplies a large share of the binding rows, so the rate says little about them. The earlier phase 3 fresh sample, drawn before the corrections, found {{p3f_err}} errors in {{p3f_verifiable}} verifiable rows ({{p3f_rate}}\%, 95\% interval {{p3f_lo}} to {{p3f_hi}}), and the phase 3 recheck of the rows phase 2 had left open found {{p3r_err}} in {{p3r_verifiable}} ({{p3r_rate}}\%). The two samples are reported separately because they were drawn at different times. The true error rate of the whole table is plausibly higher than the readable subset of the phase 4 sample suggests.

\begin{table}[t]\centering\small
{{tex_verif}}
\caption{Fresh verification sample, readable rows only.}\label{tab:vf}
\end{table}

\subsection{Completeness pilot}
A separate agent built a reference list of {{cp_ref}} candidate documents without seeing the table, and matches were made by hand. Of {{cp_pos}} candidates confirmed to contain a testing provision, {{cp_in}} are in the table, a recall of {{cp_rec}}\% (95\% interval {{cp_ci}}); a further {{cp_unres}} candidates could not be resolved either way, and counting all of them as misses gives a lower bound of {{cp_lb}}\%. The confirmed misses are all documents of the United States federal government (NIST, OMB and the national action plan). Of the {{cp_unres_miss}} unresolved candidates that are not in the table, {{cp_unres_miss_uae}} are UAE instruments, which matters because the UAE is the application focus. Matching was done by hand by the author of the table. This is a pilot under the same tool limits as the table and is not a population estimate. The table therefore should not be described as complete.

\section{The checklist tool}
The tool turns the table into a list of {{n_checklist_items}} checklist items for an organisation. It reads answers to a small set of questions (system type, decision impact, entity zone, audience, data types, acquisition route, European market, hosting) and returns for each item APPLIES, DOES\_NOT\_APPLY or UNKNOWN, with the rows that support the item. Evaluation uses a Kleene three valued logic: an answer that is missing or unknown never makes a condition true or false. For each item the tool also says which answers would settle it. Whether a source binds the organisation is reported separately and is YES only when the answers settle that the organisation is the actor the source names.

\textbf{Evaluation.} We wrote {{ev_scn}} scenarios from the checklist specification alone, with expected results from an independent oracle, and committed them before the engine was run: {{ev_spec}} with a specified result, {{ev_unspec}} invalid inputs, ten stated assumptions and five monotone pairs. Table~\ref{tab:tool} compares the engine with four naive designs.

\begin{table}[t]\centering\small
{{tex_tool}}
\caption{Item level results over {{ev_spec}} specified scenarios ({{ev_dec}} decisions). Wrongly excluded: said not applicable where applicable or unknown was expected. Wrongly applied: said applicable where not applicable or unknown was expected.}\label{tab:tool}
\end{table}

The engine agreed on {{ev_ok}} of {{ev_dec}} decisions ({{ev_acc}}\%) and matched {{ev_exact}} of {{ev_spec}} scenarios exactly, and it never wrongly excluded or wrongly applied an item. The two valued designs score about 57\% because they cannot express unknown, and they are unsafe in opposite ways: the closed design drops items the user cannot yet rule out, and the open design applies items that do not apply. The fairer comparison is with the design that ignores the answers and returns unknown for conditional items ({{ev_struct_acc}}\% with no errors), which the engine beats because it uses the answers; the oracle itself uses three valued logic, so two valued designs are penalised by construction. The scenarios were also used to find and fix two engine defects, so the figure for the final engine is after tuning to them, which is why the first run is reported. The first run, before any change, with the scenarios preregistered by commit order (the order of commits is the only timestamp, so this rests on the author's report), matched {{ev0_ok}} decisions, wrongly excluded {{ev0_wex}} items (an empty list for a multiple choice question was treated as a known answer) and crashed on {{ev0_crash}} of {{ev_invalid_n}} invalid inputs (a list given for a single value question). These defects were fixed with tests and the scenarios rerun without editing the expectations. Of the {{ev_invalid_n}} invalid inputs, {{ev_invalid_accepted}} were accepted: a comma separated string for a multiple choice question is split into a list by design, which the preregistration had not expected. One scenario still disagrees: an unknown inside a list of values, where the oracle expected APPLIES and the engine conservatively returns UNKNOWN because the specification was ambiguous. A simulation of {{ev_sim_profiles}} random profiles answered one question at a time found {{ev_sim_flips}} flips between APPLIES and DOES\_NOT\_APPLY in {{ev_sim_steps}} steps. The scenarios test that the engine implements the checklist as written; they do not show that the checklist rules are legally correct.

\textbf{Worked example.} The repository applies the tool to a synthetic, undeployed prototype system, Madoun. The answers are hypothetical and no real deployment is claimed. For this profile {{madoun_applies}} items apply, {{madoun_not}} do not and {{madoun_unknown}} remain unknown pending further answers.

\section{Threats to validity and limitations}\label{sec:limits}
\begin{itemize}
\item \textbf{Coding by language model agents.} The original coding and the second codings were produced by agents of the same family, so errors can be correlated. A human second coding and a domain or legal review have not been done.
\item \textbf{Unread sources.} {{n_blocked}} rows come from sources that could not be read and are excluded from the headline, and {{vf_unread_pct}}\% of the fresh verification sample could not be read. Error rates describe readable rows.
\item \textbf{Completeness.} The completeness pilot found a recall of {{cp_rec}}\% on confirmed candidates, with an unresolved remainder. Documents in languages other than English, and non public instruments, are under represented. Arabic originals were read for only a minority of UAE sources.
\item \textbf{Counts of provisions.} A long guidance document yields many rows and a short statute few, so counts are not league tables. The document level analysis in Section~\ref{sec:robust} partly addresses this.
\item \textbf{Currency.} Several documents were being revised at the time of retrieval. Rows carry retrieval dates, and a list of time sensitive items is kept in the repository.
\item \textbf{Not legal advice.} Applicability rules in the tool are proposals by the author and have not been reviewed by a lawyer.
\item \textbf{Related work.} All works cited here were read at abstract or landing page level, not in full text, and the literature on Gulf governance and on public sector assurance is thin in our reading list.
\end{itemize}

\section{Ethics, AI assistance and availability}\label{sec:ethics}
The study uses public documents and involves no human subjects and no personal data. Research subagents built on a large language model searched for and read sources, extracted rows and proposed codes. Separate instances verified rows and made the second codings. The author set the aims and the rules. The orchestrating agent regraded some verifier outputs against the rubric ({{p3_regraded}} in the third phase and two in the fourth), and these judgement calls are recorded in the decision log; a human review of these grades has not yet been done. Verification reduces but does not remove the risk of fabricated or misread content, and the sections above report the measured error rate. Data, code, tests and the decision log are available in the project repository (hudasol/ikhtibar); the repository records the licence terms of each part. Every number in this paper is read from a committed result file, and a test fails if a result file is older than the table it describes.

\section{Conclusion}
In this sample of documents, which were chosen because they mention testing, testing AI before deployment is mentioned far more often than it is required. In this evidence base, fewer than half of headline test provisions are binding under every assumption we tried, binding text rarely says how to test, and red teaming, content safety evaluation and human oversight validation have few binding provisions or none. The table is documented and partly independently checked, with a measured error rate and a measured gap in coverage. It is not complete and not error free. The planned next steps are a human second coding, a domain review, rereading the blocked sources, and a larger completeness benchmark.

\bibliographystyle{plainnat}
\bibliography{refs}
\end{document}
